\section*{الفصل \ref{ch:b3:hydrogeology} --- الجيولوجيا المائية}

\begin{solution}{exo:b3:hydrogeology:1}
المصفاة الأولى: $z = 120 - 40 = \qty{80}{m}$، $h = 120 - 12.5 =
\qty{107.5}{m}$، الحمل الضغطي $107.5 - 80 = \qty{27.5}{m}$. المصفاة الثانية:
$z = \qty{40}{m}$، $h = \qty{105.0}{m}$، الحمل الضغطي
\qty{65.0}{m}. الحمل أعلى عند المصفاة الضحلة، ومنه تتدفّق المياه
نحو الأسفل بتدرّج $(107.5 - 105.0)/40 = \num{0.0625}$:
تقع البيزومترات في منطقة تغذية. كان الضغط وحده سيوحي
بالعكس.
\end{solution}

\begin{solution}{exo:b3:hydrogeology:2}
$q = K\Delta h/L = 12\times0.30/0.50 = \qty{7.2}{m/d}$. المساحة هي
$\pi\times0.05^2 = \qty{7.85e-3}{m^2}$، ومنه $Q = \qty{0.0565}{m^3/d}$، أو
$56.5/24 = \qty{2.36}{L/h}$. \omterm{def:b3:hydrogeology:seepage}{سرعة التسرّب} هي $7.2/0.35 =
\qty{20.6}{m/d}$.
\end{solution}

\begin{solution}{exo:b3:hydrogeology:3}
$K = 2.5/86400 = \qty{2.894e-5}{m/s}$، و$k = K\mu/(\rho g) =
\num{2.894e-5}\times\num{1.306e-3}/(999.7\times9.81) =
\qty{3.85e-12}{m^2}$. عند \qty{20}{\celsius}، يُضرب $K$ في
$(998.2/\num{1.002e-3})/(999.7/\num{1.306e-3}) = \num{1.30}$: $K =
\qty{3.25}{m/d}$. يعود التغيّر تقريبًا كلّه إلى اللزوجة.
\end{solution}

\begin{solution}{exo:b3:hydrogeology:4}
$T = 2.5\times180 = \qty{450}{m^2/d}$. يعطي تمدّد الماء
$\num{9e-7}\times180 = \num{1.62e-4}$، نحو \qty{80}{\percent} من $S =
\num{2e-4}$؛ ويُسهم انضغاط الحجر الجيري في الباقي.
\end{solution}

\begin{solution}{exo:b3:hydrogeology:5}
حسب \cref{prop:b3:hydrogeology:thiem}، $T = Q\ln(r_2/r_1)/(2\pi(s_1 -
s_2)) = 1500\times\ln10/(2\pi\times1.53) = \qty{359}{m^2/d}$، و$K =
359/40 = \qty{9.0}{m/d}$.
\end{solution}

\begin{solution}{exo:b3:hydrogeology:6}
$Q/4\pi T = 1200/(4\pi\times300) = \qty{0.318}{m}$. بعد يوم واحد، $u =
300^2\times\num{5e-4}/(4\times300\times1) = \num{0.0375}$ و$W = -0.5772
+ 3.2834 + 0.0375 - 0.0004 = \num{2.743}$: $s = \qty{0.873}{m}$. بعد 10
أيام، $u = \num{0.00375}$، $W = -0.5772 + 5.5860 + 0.0038 = \num{5.013}$:
$s = \qty{1.596}{m}$. تكبير المدّة عشرة أضعاف يعطي أقلّ من ضعف
الانخفاض.
\end{solution}

\begin{solution}{exo:b3:hydrogeology:7}
$u \le 0.05$ عندما $t \ge r^2S/(4T\times0.05)$. عند \qty{60}{m}: $3600\times
\num{2e-4}/90 = \qty{0.008}{d}$، \qty{11.5}{min}. عند \qty{500}{m}:
\qty{0.556}{d}، \qty{800}{min}، أي أكثر من 13 ساعة. يتزايد الزمن مع $r^2$:
في اختبار يدوم يومًا أو يومين، يجب أن تكون آبار المراقبة المستخدمة في رسم
المستقيم على بُعد لا يتجاوز بضع مئات من الأمتار من بئر الضخ.
\end{solution}

\begin{solution}{exo:b3:hydrogeology:8}
في المنتصف، تبعد كلّ بئر مسافة \qty{200}{m}: $u = 200^2\times\num{2e-4}/(4
\times450\times10) = \num{4.44e-4}$، $W = \num{7.142}$، $s = 0.382\times
7.142 = \qty{2.728}{m}$ من كلّ بئر، \qty{5.456}{m} في المجموع. عند
\qty{400}{m}، $u = \num{1.78e-3}$، $W = \num{5.757}$، $s =
\qty{2.199}{m}$. تخفض كلّ بئر منسوب الماء في البئر الأخرى بمقدار
\qty{2.2}{m}: يجب أن ترفع مضخّتها الماء \qty{2.2}{m} أعلى، وهذا
يكلّف طاقة، والانخفاض المتبقّي قبل أن يجفّ مدخل المضخّة هو
أصغر.
\end{solution}

\begin{solution}{exo:b3:hydrogeology:9}
تحدّ أربعة خطوط انسياب ما بين قاعدة السدّ والقاع $N_f = 5$
قنوات؛ وتقسّم تسعة خطوط تساوي الكمون \qty{10}{m} إلى $N_d = 10$
هبوطًا كمونيًا. $q = KHN_f/N_d = 5\times10\times5/10 = \qty{25}{m^3/d}$ لكلّ متر
من السدّ، و\qty{7500}{m^3/d} لطول \qty{300}{m}. عامل الشكل المحسوب
0.512 بدلًا من 0.5، كان سيعطي \qty{25.6}{m^3/d} لكلّ متر.
\end{solution}

\begin{solution}{exo:b3:hydrogeology:10}
$k = na^2/8 = 0.30\times(\num{5e-5})^2/8 = \qty{9.4e-11}{m^2}$، و$K =
k\rho g/\mu = \num{9.375e-11}\times998.2\times9.81/\num{1.002e-3} =
\qty{9.2e-4}{m/s}$، \qty{79}{m/d}، ضمن مدى الرمل النظيف، من 0.3 إلى
\qty{150}{m/d}. يُهمل النموذج تعرّج المسارات، وهي أطول
من الأنابيب المستقيمة (وهذا يخفض $K$)، والاختناقات بين المسام،
أضيق من المسام (وهذا يخفض)، والمسام المسدودة من أحد طرفيها التي لا تحمل
تدفّقًا (وهذا يخفض)، وتباين أحجام المسام: لأنّ الأنبوب يحمل
تدفّقًا يتناسب مع $a^4$، فأوسع المسارات تحمل معظم الماء
ويرفع التباين $K$ فوق قيمة الأنابيب ذات الحجم المتوسّط.
\end{solution}

\begin{solution}{exo:b3:hydrogeology:11}
(أ)~يتلاشى اللوغاريتم عندما $2.25\,Tt = r^2S$، أي عند $r_0 =
\sqrt{2.25\,Tt/S}$. (ب)~الخزان المحصور وغير المحصور:
\[
r_0 = \sqrt{\frac{2.25\times450\times30}{\num{2e-4}}} = \qty{12.3}{km},
\qquad r_0 = \sqrt{\frac{2.25\times450\times30}{0.2}} = \qty{390}{m},
\]
$\sqrt{1000} = 31.6$ مرّة أصغر. (ج)~الضغط
ينتشر بالانتشارية $T/S$، وهي هائلة عندما $S$ صغيرة: خلال
الاختبار يصل المخروط إلى الصدوع، ومكاشف الخزان الجوفي والأنهار،
ويخرج الانخفاض عن المستقيم، وهذا ما يخبر
عالم الجيولوجيا المائية أنّه قد بلغ حدًّا.
\end{solution}

\begin{solution}{exo:b3:hydrogeology:12}
(أ)~عند النظر من الإطار المتحرّك مع الماء، تنتشر الجبهة بانحراف
معياري $\sigma_x = \sqrt{2Dt} = \sqrt{2\alpha_Lvt}$. تمرّ
$x$ تقريبًا عند $t = x/v$، حين $\sigma_x = \sqrt{2\alpha_Lx}$؛ وبالسرعة
$v$، يقابل انتشار $\sigma_x$ في المسافة انتشارًا $\sigma_t = \sigma_x/v$
في الزمن. (ب)~تقع الأزمنة \qty{16}{\percent} و\qty{84}{\percent} على بُعد
انحراف معياري واحد على جانبي المتوسّط: $\sigma_t = (247.2 -
152.6)/2 = \num{47.3}\text{ يومًا}$. ثمّ $\alpha_L = v^2\sigma_t^2/(2x) =
0.25\times47.3^2/200 = \qty{2.8}{m}$. (وقد حُسب المنحنى باستخدام
\qty{3}{m}: التقريب الطبيعي يقلّل قيمته قليلًا.) (ج)~$R =
1 + 1.9\times0.32/0.30 = \num{3.03}$: تنتقل المادة المذابة بسرعة تساوي ثلث
سرعة الماء.
\end{solution}

\begin{solution}{pb:b3:hydrogeology:1}
\textbf{1.} الطَّفل $363.2 - 282.5 = \qty{80.7}{m}$؛ الحجر الجيري $282.5 -
102.5 = \qty{180}{m}$.

\textbf{2.} $h = 363.2 - 45.0 = \qty{318.2}{m}$؛ حمل الضغط عند القمة
$318.2 - 282.5 = \qty{35.7}{m}$. يرتفع الماء \qty{35.7}{m} فوق
قمة الحجر الجيري، والطَّفل يحبسه: \omterm{def:b3:hydrogeology:aquifer-types}{الخزان الجوفي المحصور}.

\textbf{3.} من $h_2 - h_1 = 550a - 100b = -1.62$ و$h_3 - h_1 = 250a
- 600b = -1.44$: $a = \num{-2.715e-3}$ و$b = \num{1.269e-3}$ لكلّ
متر.

\textbf{4.} $|\nabla h| = \sqrt{a^2 + b^2} = \num{0.0030}$. يتدفّق الماء
على طول $-\nabla h = (\num{2.715e-3}, \num{-1.269e-3})$، شرقًا
وجنوبًا: السمت $\ang{180} - \arctan(2.715/1.269) = \ang{115}$.

\textbf{5.} النقاط قبل نحو \qty{10}{min} تقع فوق الخطّ: هناك
$u > 0.05$ ويقع منحنى ثيس فوق تقريبه اللوغاريتمي
($W > -\gamma - \ln u$، كما يبيّن شكل دالة البئر).

\textbf{6.} أفضل خطّ يمرّ بالنقاط من \qty{15}{min} فصاعدًا
\qty{1.75}{m} عند \qty{100}{min} و\qty{2.63}{m} عند \qty{1000}{min}:
$\Delta s = \qty{0.87}{m}$ لكلّ دورة لوغاريتمية.

\textbf{7.} $T = 2.3\times2160/(4\pi\times0.872) = \qty{453}{m^2/d}$.

\textbf{8.} ينخفض الخطّ \qty{0.87}{m} لكلّ دورة ابتداءً من \qty{1.75}{m} عند
\qty{100}{min}: يُبلغ الصفر قبل $1.75/0.872 = 2.01$ دورةً،
$t_0 = \qty{0.98}{min} = \qty{6.83e-4}{d}$.

\textbf{9.} $S = 2.25\,Tt_0/r^2 = 2.25\times453\times\num{6.83e-4}/3600 =
\num{1.9e-4}$، ضمن مدى الخزانات الجوفية المحصورة.

\textbf{10.} عند \qty{15}{min}، $u = 3600\times\num{1.9e-4}/(4\times453
\times0.0104) = \num{0.036}$، وأقلّ بعد ذلك.

\textbf{11.} $K = 453/180 = \qty{2.5}{m/d}$، وهو حجر جيري متشقّق.

\textbf{12.} $u = 3600\times\num{2e-4}/(4\times450\times2) =
\num{2.0e-4}$, $W = -0.5772 + 8.5172 + 0.0002 = \num{7.940}$, $s =
0.382\times7.940 = \qty{3.03}{m}$، كما سُجّل.

\textbf{13.} $u = 10^6\times\num{2e-4}/(4\times450\times30) =
\num{0.0037}$, $W = \num{5.02}$, $s = \qty{1.92}{m}$.

\textbf{14.} $r_0 = \sqrt{2.25\times450\times30/(\num{2e-4})} =
\qty{12.3}{km}$، أكبر من عرض المنطقة. بلغ المخروط
أطراف الحجر الجيري وانكشافاته والصدع، وهذا قد يجعله
يقابل صخرًا أقلّ نفاذية: لم يعد الخزان اللانهائي لثيس
يصفه، ولن يبيّن إلا سجلّ الانخفاضات ما إذا كان الحدّ
يغذّي البئر أم يوقف الانسياب.

\textbf{15.} $2160\times30 = \qty{64800}{m^3}$. ويأتي ذلك من التخزين:
الماء المنطلق بفعل تمدّد الماء وانضغاط
الحجر الجيري، عبر كامل مساحة المخروط. وبتكامل $S\,s$ على
المخروط نستعيد $Qt$.

\textbf{16.} $u = 10^6\times0.18/(4\times450\times30) = \num{3.33}$, $W
= \num{0.0086}$، $s = \qty{3}{mm}$: لن تلاحظ المزرعة ذلك.
أمّا الضخّ فسيسحب من المسام القريبة من البئر بدلًا من
ضغط الخزان الجوفي كلّه.

\textbf{17.} $q = 2.5\times0.0030 = \qty{0.0075}{m/d}$; $v = 0.0075/0.18
= \qty{0.042}{m/d}$، \qty{15}{m} سنويًا.

\textbf{18.} $200/0.0417 = \num{4800}\text{ يومًا}$، \num{13.1}~سنة؛ والملوّث،
\num{32.9}~سنة.

\textbf{19.} عبر أسطوانة نصف قطرها $r$ يعبر التصريف $Q$
مساحتها $2\pi rb$ التي تكون المسام منها نسبة $n_e$: $\mathrm dr/\mathrm
dt = -Q/(2\pi rbn_e)$، وبتكامل $2\pi bn_er\,\mathrm dr = -Q\,
\mathrm dt$ من $r$ إلى 0 نحصل على $t = \pi n_ebr^2/Q$. هنا $\pi\times0.18
\times180\times200^2/2160 = \num{1885}\text{ يومًا}$، \num{5.2}\text{ سنة} للماء
و$2.5\times5.2 = \num{12.9}\text{ سنة}$ للملوّث: البئر
يوصله أسرع بأكثر من الضعف.

\textbf{20.} $\alpha_L = \qty{20}{m}$ و$\sqrt{2\times20\times200} =
\qty{89}{m}$. يبلغ التركيز \qty{16}{\percent} من
قيمته النهائية \qty{89}{m} أمام الجبهة المتوسّطة: يصل ذلك المستوى إلى
البئر حين تكون الجبهة المتوسّطة ما زالت على بعد \qty{89}{m}، عند $(111/200)^2\times12.9 = \num{4.0}\text{ سنة}$، أي ثلث
الزمن المتوسّط. تهمّ الآثار الأولى لمياه الشرب، لا
المتوسّط.

\textbf{21.} مع: البئر يلتقط الملوّث، لذا لا يمكنه الانتشار
أبعد باتجاه التدرّج، ويمكن معالجة الماء حيث يُضخّ.
ضدّ: يسحب الضخّ الملوّث إلى بئر الإمداد أسرع ممّا
كان سيتحرّك إلى أيّ مكان، ويجب أن يُرصد الماء
ويُعالَج لعقود.

\textbf{22.} ناقلية الحجر الجيري لتلال فاردن هي $T \approx
\textbf{\qty[text-series-to-math=true]{450}{m^2/d}}$ (\qty{453}{m^2/d}
من الاختبار).
\end{solution}
